\clearpage
\VerdierPhysicalPageAnchor{0146}%
\begingroup
\small
\setlength{\parskip}{0pt}%
\setlength{\emergencystretch}{1.5em}%
\setlength{\abovedisplayskip}{3pt plus 1pt minus 1pt}%
\setlength{\belowdisplayskip}{3pt plus 1pt minus 1pt}%
\sloppy

\begingroup\itshape
\VerdierSemanticAnchor{ch2-2-3-3-f}
\noindent f) $\mathcal D$ の対象 $Z$ が
$Q^{-1}(\operatorname{Lib}_R(Q))$ に属するための必要十分条件は、右自由な対象
$Y$ と射 $s\colon Z\longrightarrow Y$ が存在し、
$s$ が $\widehat{S_{\mathcal D}(\mathcal B)}$ に属することである。ここで後者は
$S_{\mathcal D}(\mathcal B)$ の飽和化である (2.2.4, (c))。
\VerdierSemanticAnchor{ch2-2-3-3-g}
\par\noindent g) 包含関手
$L_R(Q)\subset Q^{-1}(\operatorname{Lib}_R(Q))$ は左随伴
\[
  \mathcal R\colon Q^{-1}(\operatorname{Lib}_R(Q))\longrightarrow L_R(Q)
  .
\]
をもつ。
\VerdierSemanticAnchor{ch2-2-3-3-adjunction-diagram}
\begin{center}
\begin{tikzcd}[column sep=3.2em,row sep=3.1em]
L_R(Q) \arrow[r,shift left=.7ex] \arrow[dr,"{Q|_{L_R(Q)}}"'] &
Q^{-1}(\operatorname{Lib}_R(Q))
 \arrow[l,dashed,bend right=28,"{\mathcal R}"'] \arrow[r,hook]
 \arrow[d,"{Q_1}"] & \mathcal D \arrow[d,"Q"] \\
& \operatorname{Lib}_R(Q) \arrow[ul,dashed,bend left=10,"{\Psi}"']
 \arrow[r,hook] & \mathcal D/\mathcal B
\end{tikzcd}
\end{center}
\endgroup

\VerdierSemanticAnchor{ch2-2-3-3-adjunction-explanation}
$\Psi\colon\operatorname{Lib}_R(Q)\longrightarrow L_R(Q)$ を
$Q|_{L_R(Q)}$ の擬逆関手とする（c), e) 参照）。
$\mathcal R=\Psi Q_1$ ととれる。ここで
$Q_1\colon Q^{-1}(\operatorname{Lib}_R(Q))\longrightarrow
\operatorname{Lib}_R(Q)$ は $Q$ から誘導される関手である。関手 $\Psi$ は
関手 $Q_1$ の右随伴である。

\VerdierSemanticAnchor{ch2-2-3-4}
\medskip\noindent\textbf{2.3.4. 命題 2.3.3 の証明：}\quad
\VerdierSemanticAnchor{ch2-2-3-4-a-proof}
a) を示す。まず $Y$ が $Q$-右自由ならば、関手
$X\longmapsto\operatorname{Hom}_{\mathcal D}(X,Y)$ は、
$\mathcal D/\mathcal B$ の適当な対象 $Z$ に対する関手
$X\longmapsto\operatorname{Hom}_{\mathcal D/\mathcal B}(Q(X),Z)$ と同型で
ある。よって関手 $X\longmapsto\operatorname{Hom}_{\mathcal D}(X,Y)$ は
$S_{\mathcal D}(\mathcal B)$ の射を同型に送る。これで
i) $\Rightarrow$ ii) が従う。ii) $\Rightarrow$ iii) は
$B\longrightarrow0$ が $S_{\mathcal D}(\mathcal B)$ に属することから従い、iii)
$\Rightarrow$ iv) は $S_{\mathcal D}(\mathcal B)$ の定義と (1.2.6) から
従う。iv) $\Rightarrow$ v) は次の同型 (2.2.4) から従う：
\VerdierSemanticAnchor{ch2-2-3-4-colimit-isomorphism}
\[
  \operatorname{Hom}_{\mathcal D/\mathcal B}(Q(X),Q(Y))
  \xrightarrow{\sim}
  \varinjlim_{Y\backslash S_{\mathcal D}(\mathcal B)}
  \operatorname{Hom}_{\mathcal D}(X,\mathord\cdot)\, .
\]
\VerdierSemanticAnchor{ch2-2-3-4-a-proof-complete}
最後に v) $\Rightarrow$ i) は明らかである。実際、$Y$ は $Q(Y)$ 上
右自由である。
\VerdierSemanticAnchor{ch2-2-3-4-b-c-e-proof}
b) は a), iii) から、c) は a), v) から従う。さらに a), v) により、
$\mathcal D/\mathcal B$ の右自由化可能な対象は、$\mathcal D/\mathcal B$ で
$Q$-右自由な対象の $Q$ による像と同型な $\mathcal D/\mathcal B$ の対象に
ほかならない。自由矢は同型で
ある。これで e) が従い、$\operatorname{Lib}_R(Q)$ は厳密充満三角部分圏で
あることも明らかである。
\VerdierSemanticAnchor{ch2-2-3-4-d-proof}
$Z=Z_1\oplus Z_2$ を $\operatorname{Lib}_R(Q)$ の対象とする。
$Q$-右自由な対象 $Y$ と同型 $Q(Y)\xrightarrow{\sim}Z_1\oplus Z_2$ が
存在する。したがって関手 $X\longmapsto\operatorname{Hom}_{\mathcal D}(X,Y)$
は射影子をもち、$Z$ の直和分解は $Y$ 上の射影子を定める。$Z$ の直和因子が
$\operatorname{Lib}_R(Q)$ に属するための必要十分条件はこの射影子が分解可能
であることであり、可算直和または可算積が $\mathcal D$ で表現可能ならば
そうなる (1.2.9)。
\endgroup
